\chapter{Лабораторная работа №7. Создание сети IPv6.}
\label{ch:chapter 7}

\section{Цели}

Лабораторная работа помогает получить практические навыки по изучению
следующих тем:

\begin{itemize}

\item Настройка статических адресов IPv6

\item Настройка сервера DHCPv6

\item Настройка адресов без отслеживания состояния

\item Настройка статических маршрутов IPv6

\item Способы просмотра информации IPv6

\end{itemize}

\section{Топология сети}

\includegraphics[width=0.5\textwidth]{lb7_1.png} 
\captionof{figure}{Топология сети для лабораторной работы №7}

\section{План работы}

\begin{itemize}

\item 1. Настройка статических адресов IPv6.

\item 2. Настройка сервера DHCPv6.

\item 3. Настройка назначения адресов IPv6 без отслеживания состояния.

\item 4. Вывод на экран адресов IPv6.

\end{itemize}

\section{Процедура конфигурирования}

\textbf{Шаг 1.} Настроим основные параметры устройств.

Зададим имена устройствам:

\includegraphics[width=0.5\textwidth]{lb7_2.png} 

\includegraphics[width=0.5\textwidth]{lb7_3.png} 

\includegraphics[width=0.5\textwidth]{lb7_4.png} 

\textbf{Шаг 2.} Настроим функции IPv6 на устройствах и интерфейсах.

Включим IPv6 глобально:

\includegraphics[width=0.5\textwidth]{lb7_5.png} 

\includegraphics[width=0.5\textwidth]{lb7_6.png} 

\includegraphics[width=0.5\textwidth]{lb7_7.png} 

Включим IPv6 на интерфейсе:

\includegraphics[width=0.5\textwidth]{lb7_8.png} 

\includegraphics[width=0.5\textwidth]{lb7_9.png} 

\includegraphics[width=0.5\textwidth]{lb7_10.png} 

\textbf{Шаг 3.} Настроим локальный адрес канала (link-local address) для интерфейса и проверим конфигурацию.

Настроим на интерфейсе автоматическое генерирование локального адреса канала:

\includegraphics[width=0.5\textwidth]{lb7_11.png} 

\includegraphics[width=0.5\textwidth]{lb7_12.png} 

\includegraphics[width=0.5\textwidth]{lb7_13.png} 

Выведем на экран IPv6-статус интерфейса и проверим возможность подключения:

\includegraphics[width=0.5\textwidth]{lb7_14.png} 

\includegraphics[width=0.5\textwidth]{lb7_15.png} 

\includegraphics[width=0.5\textwidth]{lb7_16.png} 

Проверим сетевое соединение между маршрутизаторами R1 и R2:

\includegraphics[width=0.5\textwidth]{lb7_17.png} 

В данном случае мы указываем интерфейс-источник или IPv6-адрес источника.

\textbf{Шаг 4.} Настроим статические IPv6-адреса на R2.

\includegraphics[width=0.5\textwidth]{lb7_18.png} 

\textbf{Шаг 5.} Настроим функцию сервера DHCPv6 на R2 и настроим R3 для получения
IPv6-адресов через DHCPv6.

Настроим функцию сервера DHCPv6:

\includegraphics[width=0.5\textwidth]{lb7_19.png} 

Настроим функцию клиента DHCPv6:

\includegraphics[width=0.5\textwidth]{lb7_20.png} 

Выведем на экран адрес клиента и информацию о DNS-сервере:

\includegraphics[width=0.5\textwidth]{lb7_21.png} 

Настроим сервер DHCPv6 для передачи адресов шлюза клиентам:

\includegraphics[width=0.5\textwidth]{lb7_22.png} 

Настроим клиент на получение маршрута по умолчанию посредством сообщение RA:

\includegraphics[width=0.5\textwidth]{lb7_23.png} 

\textbf{Шаг 6.} Настроим R1 для получения IPv6-адреса в режиме без отслеживания состояния.

Включим RA на GE 0/0/3 R2:

\includegraphics[width=0.5\textwidth]{lb7_24.png} 

\includegraphics[width=0.5\textwidth]{lb7_25.png} 

\textbf{Шаг 7.} Настроим статический маршрут IPv6.

Настроим статический маршрут на маршрутизаторе R1, чтобы обеспечить соединение между GigabitEthernet0/0/3 на маршрутизаторе R1 и GigabitEthernet0/0/3 на маршрутизаторе R3:

\includegraphics[width=0.5\textwidth]{lb7_26.png} 

Проверим возможность установления связи:

\includegraphics[width=0.5\textwidth]{lb7_27.png} 

Выведем на экран информацию о соседях IPv6:

\includegraphics[width=0.5\textwidth]{lb7_28.png} 

\textbf{Вопросы.}

\textbf{Почему интерфейс-источник должен быть указан на шаге 3 (для проверки связи между локальными адресами канала), но не должен быть указан на шаге 7 (для проверки связи между адресами GUA)?}

При проверке связи между локальными адресами канала (link-local) у маршрутизатора может быть несколько интерфейсов с адресами link-local (например, FE80::1 на разных интерфейсах). Эти адреса не уникальны в глобальном смысле — один и тот же адрес FE80::1 может существовать на разных интерфейсах. Поэтому при отправке ping на link-local адрес необходимо явно указать интерфейс-источник (через ключ \texttt{-i} или аналогичный), чтобы операционная система маршрутизатора знала, через какой интерфейс отправить запрос и какой source address использовать.

Для глобальных уникальных адресов (GUA) такой проблемы нет, поскольку GUA уникальны в сети. Маршрутизатор по таблице маршрутизации сам определяет, через какой интерфейс достижим целевой GUA, и автоматически выбирает соответствующий source address. Поэтому указывать интерфейс-источник для ping между GUA не требуется.

\textbf{Объясните, в чем отличие между конфигурацией адреса с отслеживанием состояния и конфигурацией адреса без отслеживания состояния.}

\begin{itemize}
    \item \textbf{Без отслеживания состояния (Stateless Address Autoconfiguration — SLAAC):}
    \begin{itemize}
        \item Хост самостоятельно генерирует свой IPv6-адрес на основе префикса, полученного в RA (Router Advertisement) от маршрутизатора, и своего MAC-адреса (или случайным образом).
        \item Не требуется сервер DHCPv6.
        \item Сервер не ведёт учёт (состояние) выданных адресов.
        \item Параметры, кроме адреса (например, DNS), могут отсутствовать или передаваться через опции RA (RDNSS).
    \end{itemize}
    
    \item \textbf{С отслеживанием состояния (Stateful — DHCPv6):}
    \begin{itemize}
        \item Адрес (и другие параметры, такие как DNS-серверы, домен) выдаётся централизованно сервером DHCPv6.
        \item Сервер ведёт базу данных (состояние) выданных адресов и их привязку к клиентам.
        \item Позволяет централизованное управление, учёт, продление аренды, принудительное освобождение адресов.
        \item Используется флаг \texttt{M} (Managed) в RA, чтобы указать клиентам на использование DHCPv6.
    \end{itemize}
\end{itemize}

Также возможна комбинированная схема: адрес через SLAAC, а дополнительные параметры (DNS и т.д.) — через DHCPv6 (флаг \texttt{O} — Other Configuration).


